%قبل از شروع مسئله ترجیحاً چیزی ننویسید. اما اگر خواستید بنویسید، در قالب لیستی (فقرات و شمارش) بنویسید.
%برای اضافه کردن نمره، بعدی از شروع‌مسئله، [x] را اضافه کنید و به جای x  نمره مورد نظر را قرار دهید. برای زیرمسئله‌ها هم به همین صورت است.
%\شروع‌مسئله[10]
\شروع‌مسئله [25]
% صورت اصلی مسئله را در این‌جا بنویسید.
گادوین مِنشا،  به تازگی از تیم سابقش کناره‌گیری کرده و برای کار کردن در یک مؤسسۀ آموزشی استخدام شده است. وظیفۀ او در مؤسسه به این شرح است: در روز به او تعدادی درخت دودویی داده می‌شود و او باید با  جابه‌جایی‌ تعدادی از مقادیر رأس‌های درخت کاری کند که این درخت دودویی «بارِ دانش» بگیرد. درخت‌هایی که ‌بارِ دانش گرفته‌اند، در حقیقت درخت دودویی جستجو هستند.به این دلیل که منشا دنبال  ترفیع مقامش است، تصمیم گرفته برنامه‌ای بنویسد که  این کار را انجام دهد تا بتواند نشان «چرخ نیلوفری» را از رئیسش  بگیرد. از آن‌جا که منشا اخیراً به دانشگاه شریف آمده و دوستانی پیدا کرده‌است، از آن‌ها خواسته برایش برنامه‌ای بنویسند تا بتواند نشان «چرخ نیلوفری»  را دریافت کند.

\متن‌سیاه{ورودی:}

در خط اول ورودی ابتدا $n$ می‌آید که نشان‌دهندۀ تعداد رأس‌های درخت دودویی اولیه است.
در خط دوم ورودی $n$ عدد می‌آیند که نشان‌دهندۀ مقدار  رأس‌های درخت دودویی هستند.
نحوۀ ورودی دادن درخت به این‌ترتیب است که اگر ورودی $2i$اُم و $2i+1$اُم به ترتیب فرزند چپ و راست ورودی $i$اُم هستند.

\متن‌سیاه{خروجی:}

در تنها خط خروجی، کم‌ترین تعداد جابه‌جایی‌ها برای تبدیل درخت دودویی به ددج را چاپ کنید.

\متن‌سیاه{محدودیت‌ها:}

تمام اعداد ورودی از $10^5$ کم‌تر هستند.

\متن‌سیاه{ورودی و خروجی نمونه:}

% تست اول
\begin{latin}
	\begin{tabularx}{\textwidth}{|X | X|}
		\hline
		\textbf{\rl{ورودی }} &
		{\centering\textbf{\rl{خروجی }}}\\
		\hline
		3&
		1\\
		1 2 3&\\
		\hline
		
	\end{tabularx}
\end{latin}

\شروع{شکل}[h!]
\تنظیم‌ازوسط
  \درج‌تصویر[پهنا=0.5\پهنای‌متن]{src/fig.png}
\پایان{شکل}

%اگر مسئله چند قسمت داشت به این صورت عمل کنید:
%\شروع{زیرمسائل}
%\زیرمسئله[0] % مسئلۀ ۱آ
%\زیرمسئله % مسئلۀ ۱ب
%\زیرمسئله % مسئلۀ ۱پ
%\زیرمسئله % مسئلۀ ۱ت
%\پایان{زیرمسائل}

\solution{
	پاسخ نمونه.
}




%%%%%%%%%%%%%%%%%%%%%%%%%%%%%%%%%%%%%%%%%%%%%%%%%%%%%
% برخی دستورات مفید برای تمرین‌ها. %
%%%%%%%%%%%%%%%%%%%%%%%%%%%%%%%%%%%%%%%%%%%%%%%%%%%%%

% فرمول‌های نمونه:
%$f(x) = \Theta(x)$
%$T(x, y) \leq \log(x) + 2^y + \binom{2n}{n}$
% {\tt code\_function}


% لیست موارد بدون شماره‌گذاری:
%\شروع{فقرات}
%    \فقره مورد اول 
%        \شروع{فقرات}
%            \فقره مورد اول داخلی 
%                \شروع{فقرات}
%                    \فقره اولین مورد خیلی داخلی 
%                    \فقره دومین مورد خیلی داخلی 
%                \پایان{فقرات}
%            \فقره مورد دوم داخلی 
%        \پایان{فقرات}
%    \فقره مورد دوم 
%\پایان{فقرات}

% لیست موارد شماره‌گذاری شده:
%\شروع{شمارش}
%    \فقره یک مورد 
%    \فقره یک مورد دیگر 
%    \فقره باز هم یک مورد دیگر
%\پایان{شمارش}

% برای نوشتن یک فرمول چند خطی که یک شماره بخوره. برای حالت‌های دیگر به سایت زیر مراجعه کنید:
% https://www.overleaf.com/learn/latex/Aligning_equations_with_amsmath
%\begin{align} 
%    \begin{split}
%        (x+y)^3 &= (x+y)^2(x+y) \\
%                &= (x^2+2xy+y^2)(x+y) \\
%                &= (x^3+2x^2y+xy^2) + (x^2y+2xy^2+y^3) \\
%                &= x^3+3x^2y+3xy^2+y^3
%    \end{split}                                 
%\end{align}

% نوشتن ماتریس و آرایه دوبعدی
%\begin{align}
%    A = 
%    \begin{bmatrix}
%        A_{11} & A_{21} \\
%        A_{21} & A_{22}
%    \end{bmatrix}
%\end{align}

% اضافه کردن تصویر با شرح. برای اضافه کردن تصویر بدون شرح، فقط دستور درج‌تصویر را استفاده کنید.
%\شروع{شکل}[h!]
%\تنظیم‌ازوسط
%  \درج‌تصویر[پهنا=0.5\پهنای‌متن]{img.jpg}
%\شرح{شرح تصویر}
%\برچسب{تصویر:۱}
%\پایان{شکل}
%\پایان{زیرمسائل}

%کد زیر را ببینید.
%%\lstinputlisting{source_filename.py}
%\begin{latin}
%\begin{lstlisting}
%# Import the modules
%import sys
%import random
%
%ans = True
%
%while ans:
%question = raw_input("Ask the magic 8 ball a question: (press enter to quit) ")
%
%answers = random.randint(1,8)
%
%if question == "":
%sys.exit()
%
%elif answers == 1:
%print "It is certain"
%
%elif answers == 2:
%print "Outlook good"
%
%elif answers == 3:
%print "You may rely on it"
%
%elif answers == 4:
%print "Ask again later"
%
%elif answers == 5:
%print "Concentrate and ask again"
%
%elif answers == 6:
%print "Reply hazy, try again"
%
%elif answers == 7:
%print "My reply is no"
%
%elif answers == 8:
%print "My sources say no"
%\end{lstlisting}
%\end{latin}
